% 379_Lien_Plenum_Photon_De_ch.tex
% Johann Pascher, 29. Sept. 2026
% Das akustische Plenum und das Photon: Lien (2026) im Licht von
% Dok. 267, 290, 340, A265 und R128

\providecommand{\BM}{\textbf{[B]}}
\providecommand{\KM}{\textbf{[K]}}
\providecommand{\SM}{\textbf{[S]}}
\providecommand{\XM}{\textbf{[X]}}
\providecommand{\QM}{\textbf{[Q]}}

\chapter*{Das akustische Plenum und das Photon}
\addcontentsline{toc}{chapter}{Das akustische Plenum und das Photon}

\begin{abstract}
C.~C.~Lien (Zenodo, DOI 10.5281/zenodo.21800087, 8.~September 2026) deutet den
Raum als kompressibles Superfluid, leitet die Elektronmasse aus der
Temperatur des kosmischen Mikrowellenhintergrunds ab, übernimmt für Myon und
Tau eine Massenformel, setzt die Neutrinomasse mit $k_BT_\text{CMB}$ gleich und
beschreibt das Photon als longitudinale Schockwelle mit transversaler
Nachwelle. Dieses Dokument prüft die Arbeit als Fallstudie mit einem
Prüfskript (20/20). Die Zahlenangaben sind reproduzierbar \KM. Keine der
tragenden Aussagen hält der Prüfung stand: der Elektronentreffer liegt
88~Standardabweichungen neben dem Messwert und ist im Nulltest nicht
signifikant, die Leptonformel ist bekannt, die Neutrinomasse ist mit der
Oszillation unverträglich, die Rotverschiebung wäre frequenzabhängig, und
das Photon hätte drei statt zwei Polarisationen \XM. Die Gegenprobe zeigt,
wie die FFGFT an denselben Stellen verfährt, ohne deren Inhalte zu
wiederholen.
\end{abstract}

% ============================================================
\section*{Statusmarkierungen}
% ============================================================
\begin{center}
\begin{tabular}{@{}lp{10cm}@{}}
\textbf{[B]} & Algebraisch bewiesen. \\[3pt]
\textbf{[K]} & Numerisch bestätigt (Prüfskript). \\[3pt]
\textbf{[S]} & Strukturell plausibel, offen. \\[3pt]
\textbf{[X]} & Geprüft und verworfen. \\[3pt]
\textbf{[Q]} & Aus der Quelle (Lien 2026) wörtlich übernommen. \\
\end{tabular}
\end{center}

% ============================================================
\section{Die Quelle}
% ============================================================

Lien ersetzt den leeren Raum durch ein Medium:
\begin{quote}\QM
\enquote{By treating the spatial manifold as a continuous compressible
superfluid, we derive the foundational particle phase-states (electron, muon,
tau, and neutrino) from the Cosmic Microwave Background (CMB) acoustic
baseline.} (Conclusion)
\end{quote}
Der Hintergrund ist stationär, die Rotverschiebung keine Expansion:
\begin{quote}\QM
\enquote{Cosmological redshift is therefore not temporal expansion; it is
Acoustic Pulse Broadening -- the observation that waves of higher energetic
frequencies decay faster than those of a lower energetic frequency when
traversing the same medium.} (Abschnitt II.A)
\end{quote}
Das Photon ist kein Teilchen:
\begin{quote}\QM
\enquote{Photonic emission is therefore not a discrete point-particle
phenomenon. It is a primary longitudinal acoustic shockwave [\ldots] which
structurally interacts with the azimuthal shear of its source to continuously
generate an orthogonal transverse spiral wake.} (Abschnitt V)
\end{quote}
Die zentralen Formeln sind
\begin{align}
f_e &= 2\pi\,(\alpha^{-1})^4\,\frac{k_BT_\text{CMB}}{h}, \label{eq:lien-fe}\\
m_n &= m_e\Bigl(1+\tfrac{3}{2}\,\alpha^{-1}\sum_{k=0}^{n}k^4\Bigr), \label{eq:lien-mn}\\
m_\nu &= k_BT_\text{CMB}/c^2. \label{eq:lien-mnu}
\end{align}

% ============================================================
\section{Nachrechnung}
% ============================================================

Alle sechs Zahlenangaben sind reproduzierbar \KM:
$f_\text{CMB}=5{,}679\times10^{10}$\,Hz, $f_e=1{,}2583\times10^{20}$\,Hz
gegen die Compton-Frequenz $1{,}2356\times10^{20}$\,Hz ($+1{,}84\,\%$),
$m_\mu/m_e=206{,}55$ ($-0{,}10\,\%$), $m_\tau/m_e=3495{,}4$ ($+0{,}52\,\%$) und
$m_\nu=0{,}235$\,meV. Rechenfehler liegen nicht vor. Die Frage ist allein,
was die Zahlen bedeuten.

% ============================================================
\section{Signifikanz}
% ============================================================

\paragraph{Der Elektronentreffer ist ein Fehlschlag.}
Das Verhältnis $m_ec^2/(k_BT_\text{CMB})=2{,}176\times10^9$ liegt $1{,}84\,\%$
unter $2\pi\alpha^{-4}=2{,}216\times10^9$. $T_\text{CMB}$ ist auf $0{,}02\,\%$
bekannt; die Abweichung beträgt damit rund 88~Standardabweichungen \KM. Lien
deutet sie als \enquote{dynamic acoustic drag}, berechnet diesen Widerstand
aber nicht. Eine nachträglich benannte Reibung ist keine Korrektur, sondern
ein zweiter freier Parameter.

\paragraph{Nulltest.}
Die Formelfamilie $a\,\pi^b\alpha^{-k}$ mit dreizehn einfachen Brüchen $a$,
$b\in\{-3,\dots,3\}$ und $k\in\{1,\dots,6\}$ umfasst 546~Formeln. Gegen
zufällige Zielwerte zwischen $10^9$ und $10^{10}$ gelingt ein Treffer auf
$1{,}84\,\%$ in 61\,\% der Fälle \KM. Der Treffer ist daher nicht signifikant.
Die FFGFT wendet denselben Maßstab auf ihre eigenen Rasterbefunde an
(Dok.~342, 343 §G$'$): Was der Zufall ebenso liefert, wird nicht gewertet.

\paragraph{Die Leptonformel ist bekannt.}
Gl.~\eqref{eq:lien-mn} ist die Formel von A.~O.~Barut (Phys.\ Rev.\ Lett.\
42, 1251, 1979) ohne Änderung \BM. Die Genauigkeit für Myon und Tau ist
Baruts Ergebnis, nicht das des Plenums. Für $n=3$ ergäbe dieselbe Formel ein
geladenes Lepton bei $10{,}3$\,GeV \KM, das experimentell ausgeschlossen ist;
Liens Begründung über die Chirikov-Überlappung bleibt qualitativ. Zur
Generationenzahl in der FFGFT siehe Dok.~317 und 316.

% ============================================================
\section{Rotverschiebung}
% ============================================================

\paragraph{Dämpfung ändert die Amplitude, nicht die Frequenz.}
Stokes-Dämpfung $\alpha\propto f^2$ verringert die Amplitude einer Welle; ihre
Frequenz bleibt erhalten. Numerisch behält eine gedämpfte 50-Hz-Schwingung
ihr Spektralmaximum bei 50\,Hz \KM. Eine Rotverschiebung entsteht so nicht.

\paragraph{Eine frequenzabhängige Rotverschiebung ist ausgeschlossen.}
Wäre $z$ an eine Dämpfung $\propto f^2$ gebunden, unterschieden sich die
21-cm-Linie und Lyman-$\alpha$ um einen Faktor $3\times10^{12}$ \KM. Gemessen
stimmt $z$ zwischen beiden auf etwa $10^{-5}$ überein. Die FFGFT hat diesen
Fehler selbst durchlaufen: Ältere Korpusdokumente führten eine chromatische
Rotverschiebung, die A265 als überholt ausweist.

\paragraph{Lichtkurvendehnung.}
Lien erwähnt die $(1+z)$-Dehnung der Supernova-Lichtkurven nur als Argument
der Standardkosmologie, ohne sie zu erklären. Ein statisches Bild muss sie
enthalten. In der FFGFT folgt sie aus $\tilde T\cdot m=1$: Wellenlänge und
Uhrentakt werden gemeinsam gebucht, die Dehnung ist für alle $z$ exakt
$(1+z)$ (Dok.~312, A265, R128) \BM. Nach Dok.~267 sind alle Frequenz- und
Taktsignaturen zwischen statischem Bild und Expansion entartet; ein
statischer Hintergrund ist damit zulässig, aber nur, wenn er achromatisch
ist. Liens Plenum ist es nicht \XM.

% ============================================================
\section{Das Photon}
% ============================================================

Licht hat im Vakuum genau zwei transversale Polarisationen und keine
longitudinale Komponente. Ein masseloses $U(1)$-Eichfeld hat vier
Komponenten, von denen Eichfreiheit und Zwangsbedingung zwei entfernen
\BM. Liens Bild aus longitudinaler Primärwelle und transversaler Nachwelle
hätte drei Freiheitsgrade \XM. Das barokline Drehmoment
$\nabla\rho\times\nabla P$ erzeugt eine Wirbelkomponente in $z$-Richtung,
aber keine Welle mit zwei Polarisationszuständen. Photonenzählung
(Photoeffekt, Antibunching) behandelt die Arbeit nicht.

Das Argument, ein Teilchenbild müsse im Fernfeld zu einer
\enquote{quantisierten Streuung} zerfallen, verwechselt die Wellenfront
eines klassischen Feldes mit der Zählstatistik seiner Quanten. In der FFGFT
ist das Photon exakt masselos, der reine Phasen- und Verbindungsanker, und
$E_\text{bit}=\hbar c/L$ ist eine Photonenenergie (Dok.~290). Es braucht dafür
kein Medium.

% ============================================================
\section{Neutrino}
% ============================================================

Eine einzige Masse von $0{,}235$\,meV kann die atmosphärische
Massendifferenz $\Delta m^2_\text{atm}\approx2{,}5\times10^{-3}\,\text{eV}^2$
nicht erzeugen. Dafür muss mindestens ein Zustand etwa 50\,meV schwer sein
\KM. Die Gleichsetzung mit $k_BT_\text{CMB}$ ist damit ausgeschlossen \XM. Die
FFGFT-Hierarchie aus Dok.~340 erfüllt diese Bedingung mit ihrem schwersten
Zustand bei $49{,}96$\,meV \KM.

% ============================================================
\section{Der gemeinsame Anker}
% ============================================================

Auffällig ist, dass Lien dieselbe Zahl berührt, die in der FFGFT der
SI-Anker des kosmischen Sektors ist: $m_ec^2/(k_BT_\text{CMB})=2{,}176\times10^9$
(A265). Die FFGFT führt dieses Verhältnis ausdrücklich als Anker, nicht als
Herleitung; der Exponent $41/4$ der Brückenformel ist offen (P20), und der
kosmische Sektor ist bewusst ausgeklammert (P39). Liens Kandidat
$2\pi\alpha^{-4}$ verfehlt die Zahl um $1{,}84\,\%$ und ist im Nulltest nicht
von Zufall zu unterscheiden. Er schließt die Lücke daher nicht; er zeigt,
wie leicht sich an dieser Stelle Näherungen finden lassen \SM.

% ============================================================
\section{Was übereinstimmt und was nicht}
% ============================================================

Zwei Ausgangspunkte teilt die Arbeit mit der FFGFT: einen stationären
statt expandierenden Hintergrund und die Ablehnung punktförmiger
Singularitäten. In der FFGFT sichert die kleinste Länge
$L_0=\xi\,\ell_P$ die Endlichkeit (Dok.~180, R90); bei Lien leistet das ein
Dichtemaximum $\rho_E$. Dieses ist jedoch zirkulär, weil der klassische
Elektronenradius $r_e=\alpha\hbar/(m_ec)$ die Elektronmasse bereits enthält
\KM.

Der Unterschied liegt in der Methode. Lien führt ein Medium mit
Materialkonstanten ein (Dämpfung, Viskosität, Widerstand) und erklärt
Abweichungen nachträglich durch diese Konstanten. Die FFGFT setzt eine
kompakte Geometrie mit einem Parameter $\xi$, leitet Zahlen ab und prüft sie
mit festen Kriterien. Wo eine Abweichung bleibt, wird sie als offene Brücke
geführt, nicht wegerklärt.

% ============================================================
\section{Zusammenfassung}
% ============================================================

\begin{center}
\begin{tabular}{>{\raggedright\arraybackslash}p{4.2cm}>{\raggedright\arraybackslash}p{4.8cm}>{\raggedright\arraybackslash}p{2.2cm}l}
\toprule
Aussage (Lien) & Prüfbefund & FFGFT & Status \\
\midrule
$f_e=2\pi\alpha^{-4}f_\text{CMB}$ & $+1{,}84\,\%$, 88\,$\sigma$, Nulltest 61\,\% & A265, P20 & \XM \\
Leptonformel Gl.~\eqref{eq:lien-mn} & Barut 1979, $n=3$ ausgeschlossen & 316, 317 & \XM \\
$m_\nu=k_BT_\text{CMB}$ & unverträglich mit $\Delta m^2_\text{atm}$ & 340 & \XM \\
Rotverschiebung durch Dämpfung & frequenzabhängig, ändert $f$ nicht & 267, 312, R128 & \XM \\
Photon longitudinal + transversal & drei statt zwei Freiheitsgrade & 290 & \XM \\
Stationärer Hintergrund & zulässig, wenn achromatisch & 267, R128 & \SM \\
\bottomrule
\end{tabular}
\end{center}

Die Arbeit verbindet bekannte Formeln, eine nachträglich benannte Reibung
und ein anschauliches Medium zu einem geschlossen wirkenden Bild. Die
Nachrechnung zeigt, dass die Zahlen stimmen, die Deutung aber an jeder
tragenden Stelle an Messdaten scheitert. Der brauchbare Kern, ein statischer
Hintergrund ohne Singularitäten, ist in der FFGFT bereits anders und
achromatisch umgesetzt.

\paragraph{Prüfskript.} \texttt{2/python/Dok379\_Skripte/pruef\_379\_lien\_photon.py}
(20/20): Nachrechnung (A), Signifikanz und Nulltest (B), Messdaten (C),
Gegenprobe FFGFT (D). Zwei Eingänge sind Literaturwerte und als solche
gekennzeichnet: die Übereinstimmung von $z$ zwischen 21~cm und Lyman-$\alpha$
und der DES-Wert $b=1{,}003\pm0{,}011$.

\paragraph{Quelldokument.} Lien, C.~C.: \textit{Electromagnetic Illusion:
Photonic Emission as Orthogonal Acoustic Shockwaves in a Compressible
Superfluid Plenum.} Zenodo (2026), DOI 10.5281/zenodo.21800087.
Barut, A.~O.: \textit{Lepton Mass Formula.} Phys.\ Rev.\ Lett.\
\textbf{42}, 1251 (1979).
